\documentclass[11pt]{article}
\usepackage{amsmath,amssymb}
\usepackage[margin=1in]{geometry}
\allowdisplaybreaks
\begin{document}

\section*{Two-step-ahead hybrid adaptive predictive control}

\begin{equation}
\operatorname*{Min}_{S(k)} J
=\sum_{j=1}^{F}\left[
\sum_{h=1}^{H(k+2)}p_h^{\Delta T(k+2)}(k+2)\cdot
\left.C_j(k+2)\right|_{S_j(k),h}
-\overbrace{C_j(k)}^{\text{known constant}}
\right].
\tag{15, final line}
\end{equation}

\paragraph{Subject to the source-defined routing and prediction conditions.}
The sequence decisions obey the insertion and precedence conditions below, and their predicted loads and departure times are generated by (6)--(9). The paper does not print a separate algebraic subject-to block beside (15). Its capacity treatment is stated in the solution procedure: a sequence infeasible for capacity is not repaired, but receives a penalized fitness. Neither a general capacity inequality nor a formula for that penalty is supplied there; neither is invented here.

\section*{Fleet, control sequences, and admissible insertions}
The service region is an influence area $A$, with a service network of length $D$. The vehicle set is $V$, of size $F$. Service requests arrive in real time, at a rate denoted by $\mu$. Vehicle indices are $j\in V$, written $j=1,\ldots,F$ in the formulas. At an event instant $k$, the known previous sequence set is $S(k-1)$, and the routing decision is
\begin{equation}
u(k)=S(k)=\{S_1(k),\ldots,S_j(k),\ldots,S_F(k)\}.
\tag{3}
\end{equation}
For an individual vehicle,
\[
u_j(k)=S_j(k)
=\bigl[s_j^1(k)\ \cdots\ s_j^i(k)\ \cdots\ s_j^{w_j(k)}(k)\bigr]^T,
\]
where $s_j^i(k)$ represents the $i$th stop and $w_j(k)$ is the number of stops. The detailed representation of that sequence is
\begin{equation}
S_j(k)=
\begin{bmatrix}
r_j^1(k)&1-r_j^1(k)&\Gamma_j^1(k)&\mathrm{label}_j^1(k)\\
\vdots&\vdots&\vdots&\vdots\\
r_j^i(k)&1-r_j^i(k)&\Gamma_j^i(k)&\mathrm{label}_j^i(k)\\
\vdots&\vdots&\vdots&\vdots\\
r_j^{w_j(k)}(k)&1-r_j^{w_j(k)}(k)&\Gamma_j^{w_j(k)}(k)&\mathrm{label}_j^{w_j(k)}(k)
\end{bmatrix}_{w_j(k)\times4}.
\tag{7}
\end{equation}
For $i=1,\ldots,w_j(k)$,
\[
r_j^i(k)=
\begin{cases}
1,&\text{if stop $i$ belonging to $S_j(k)$ is a pickup},\\
0,&\text{otherwise}.
\end{cases}
\]
The first two columns identify a pickup by $[1\ 0]$ and a delivery by $[0\ 1]$. The external travel-time quantity $\Gamma_j^i(k)$ is the total travel time between points $i-1$ and $i$, plus the transfer-operation delay at node $i$. The passenger identifier $\mathrm{label}_j^i(k)$ is used to check precedence: the pickup must occur before the delivery of the same client. No particular analytic travel-time function is specified in this model. Constant vehicle speed is an assumption used for simulation, not an additional general constraint imposed here.

The selected insertion design does not swap previously scheduled stops. The new request contributes a pickup and a delivery to the sequence of the vehicle receiving it. The paper describes the resulting matrix as containing the $w_j(k-1)$ previously scheduled rows and the two new rows. Insertion positions are denoted by
\[
\vartheta=(pu,de),
\qquad
pu\in\{1,\ldots,w_j(k)-1\},
\qquad
de\in\{2,\ldots,w_j(k)\},
\qquad pu<de.
\]
Here $pu$ and $de$ are integer row positions, not the continuous particles used by the solver. The $pu$th and $de$th rows of the new sequence are the pickup and delivery of the incoming request. These conditions apply to insertions for the current call and to insertions of each predicted call into the corresponding future sequence. The explicit feasible-pair test in the paper's procedure is
\[
\text{If }pu_l<de_l,\quad\text{then }\vartheta_l=(pu_l,de_l).
\]
The subscript $l$ in this last source relation indexes a candidate insertion; it is not a passenger or vehicle index. Particle generation, rounding, repair, and particle-motion equations are solver operations, not additional constraints of (15).

Vehicles follow $S(k-1)$ until the new call triggers a routing decision. A chosen $S_j(k)$ remains fixed during $(k,k+1)$ except that reaching a scheduled pickup or delivery removes the accomplished stop. The source describes deletion of the first sequence row when its task is satisfied. Thus insertion expands, and completed service reduces, the sequence and state dimensions. The two-row expansion concerns the vehicle receiving the new request, rather than adding that request to every vehicle.

\section*{State variables and predictive dynamics}
At instant $k$, the state refers to $S(k-1)$, before insertion of the new call. $L_j^i(k)$ and $T_j^i(k)$ are, respectively, the vehicle load and clock time of departure after vehicle $j$ leaves stop $i$, computed at instant $k$. The zero-index quantities refer to the current vehicle position, as depicted in Figure 2. For every $j\in V$,
\begin{equation}
L_j(k)=
\bigl[L_j^0(k)\ L_j^1(k)\ \cdots\ L_j^{w_j(k-1)}(k)\bigr]^T_{(w_j(k-1)+1)\times1},
\tag{4}
\end{equation}
\begin{equation}
T_j(k)=
\bigl[T_j^0(k)\ T_j^1(k)\ \cdots\ T_j^{w_j(k-1)}(k)\bigr]^T_{(w_j(k-1)+1)\times1}.
\tag{5}
\end{equation}
The system-wide collections and variable roles are
\[
\begin{aligned}
L(k)&=\{L_1(k),\ldots,L_j(k),\ldots,L_F(k)\},\\
T(k)&=\{T_1(k),\ldots,T_j(k),\ldots,T_F(k)\},\\
x(k)&=\{L(k),T(k)\},\qquad y(k)=T(k),\qquad u(k)=S(k).
\end{aligned}
\]
Departure times are continuous state variables; loads describe discrete passenger numbers. The stochastic load and departure-time variables are replaced in the prediction by their conditional expectations, not by additional freely chosen state decisions:
\begin{equation}
\begin{aligned}
\hat{x}_j(k+1)
&=\begin{bmatrix}
E\{L_j(k+1)/k\}\\ E\{T_j(k+1)/k\}
\end{bmatrix}
=\begin{bmatrix}\hat{L}_j(k+1)\\\widehat{T}_j(k+1)\end{bmatrix}\\
&=\begin{bmatrix}
f_L(L_j(k),S_j(k))\\
f_T(T_j(k),S_j(k))
\end{bmatrix},\qquad \forall j=1,\ldots,F.
\end{aligned}
\tag{6}
\end{equation}
The slash in $E\{\cdot/k\}$ retains the paper's conditional-expectation notation. The prediction models are
\begin{equation}
\hat{L}_j(k+1)=f_L(L_j(k),S_j(k))
=A_LL_j(k)+B_L(S_j(k)),
\tag{8}
\end{equation}
with the following complete matrices:
\[
B_L(S_j(k))=B_L^2\cdot(S_j(k)\cdot B_L^1),
\]
\[
A_L=
\begin{bmatrix}
1&0&\cdots&0\\
1&0&\cdots&0\\
\vdots&\vdots&\ddots&\vdots\\
1&0&\cdots&0
\end{bmatrix}_{(w_j(k)+1)\times(w_j(k-1)+1)},
\qquad
B_L^1=\begin{bmatrix}1\\-1\\0\\0\end{bmatrix},
\]
\[
B_L^2=
\begin{bmatrix}
0&0&0&\cdots&0&0\\
1&0&0&\cdots&0&0\\
1&1&0&\cdots&0&0\\
1&1&1&\ddots&\vdots&\vdots\\
\vdots&\vdots&\vdots&\ddots&0&0\\
\vdots&\vdots&\vdots&&1&0\\
1&1&1&\cdots&1&1
\end{bmatrix}_{(w_j(k)+1)\times w_j(k)}.
\]
The departure-time model is
\begin{equation}
\widehat{T}_j(k+1)=f_T(T_j(k),S_j(k))
=A_T\cdot T_j(k)+B_T(S_j(k)),
\tag{9}
\end{equation}
where
\[
B_T(S_j(k))=B_T^2\cdot(S_j(k)\cdot B_T^1),
\]
\[
A_T=
\begin{bmatrix}
1&0&\cdots&0\\
1&0&\cdots&0\\
\vdots&\vdots&\ddots&\vdots\\
1&0&\cdots&0
\end{bmatrix}_{(w_j(k)+1)\times(w_j(k-1)+1)},
\qquad
B_T^1=\begin{bmatrix}0\\0\\1\\0\end{bmatrix},
\]
\[
B_T^2=
\begin{bmatrix}
0&0&0&\cdots&0&0\\
1&0&0&\cdots&0&0\\
1&1&0&\cdots&0&0\\
1&1&1&\ddots&\vdots&\vdots\\
\vdots&\vdots&\vdots&\ddots&0&0\\
\vdots&\vdots&\vdots&&1&0\\
1&1&1&\cdots&1&1
\end{bmatrix}_{(w_j(k)+1)\times w_j(k)}.
\]
The matrices are used for each vehicle with its own sequence lengths. Their sizes adapt as stops are inserted or completed. The same prediction models are used over the second prediction step with the future sequences. No additional event-time transition formula is supplied for advancing the current-position quantities between calls.

\section*{Vehicle costs and the selected two-step horizon}
The selected prediction horizon is $N=2$. The controller computes the current and future sequence decisions and applies only $S(k)$, following the receding-horizon rule. Future sequences $S_j(k+1)$ depend on the predicted request $h$ and the preceding sequence $S_j(k)$; the conditioning bars in the cost expressions retain this dependence without renaming the source variables.

The vehicle-cost definition inherited from the parent objective is
\begin{equation}
\begin{aligned}
\left.C_j(k+t)\right|_{S_j(k+t-2),h}
={}&\left.\sum_{i=1}^{w_j(k+t-1)}\Bigl\{
\underbrace{[\hat{L}_j^{i-1}(k+t)+1]
\bigl(\widehat{T}_j^i(k+t)-\widehat{T}_j^{i-1}(k+t)\bigr)}_{J_{\mathrm{travel\ time}}}\\
&\hspace{16mm}+\underbrace{r_j^i(k+t-1)\alpha
\bigl(\widehat{T}_j^i(k+t)-T_j^0(k+t)\bigr)}_{J_{\mathrm{waiting\ time}}}
\Bigr\}\right|_{S_j(k+t-2),h}.
\end{aligned}
\tag{11}
\end{equation}
Here $j=1,\ldots,F$, and for the selected horizon $t=1,2$ and $h=1,\ldots,H(k+t)$. The printed second-step specialization is
\begin{equation}
\begin{aligned}
\left.C_j(k+2)\right|_{S_j(k),h}
={}&\left.\sum_{i=1}^{w_j(k+1)}\Bigl\{
\underbrace{[\hat{L}_j^{i-1}(k+2)+1]
\bigl(\widehat{T}_j^i(k+2)-\widehat{T}_j^{i-1}(k+2)\bigr)}_{J_{\mathrm{travel\ time}}}\\
&\hspace{16mm}+\underbrace{r_j^i(k+1)\alpha
\bigl(\widehat{T}_j^i(k+2)-T_j^0(k+2)\bigr)}_{J_{\mathrm{waiting\ time}}}
\Bigr\}\right|_{S_j(k),h}.
\end{aligned}
\tag{16}
\end{equation}
The expected load $\hat{L}_j^{i-1}(k+t)$ is the passenger load on segment $(i-1,i)$. The difference of consecutive expected departure times includes travel along that segment and transfer delay at node $i$. The added $1$ in the travel-time term accounts for the driver and serves as a proxy for operating cost, including empty travel. The waiting-time difference is measured from the current-position time $T_j^0(k+t)$, which is printed without a hat. The coefficient $\alpha$ weights waiting time; the general formulation supplies no numerical value or sign restriction for it. The simulation choice $\alpha=1$ is not imposed here.

\section*{Request probabilities and prediction time}
The event step has variable duration: it is the interval between consecutive calls. The source defines $k+t$ as the instant at which the $t$th request enters, measured from time interval $k$. The notation $\Delta T(k+t)$ identifies the time interval to which event step $k+t$ belongs; it is not defined as an algebraic difference of the vehicle departure-time vectors.

$H(k+t)$ is the number of probable request types at that instant. A type $h$ is an origin--destination pair of clusters, with
\[
h\in\{1,\ldots,H(k+t)\}.
\]
The current new request is known. The first-step and second-step probability conditions used in (15) are
\[
H(k+1)=1,\qquad p_1^{\Delta T(k+1)}(k+1)=1,
\qquad
\sum_{h=1}^{H(k+2)}p_h^{\Delta T(k+2)}(k+2)=1.
\]
For the historical-data version used in this paper, the study area is partitioned into clusters. The call mass center is the geographical location of the most likely call during a specified period and within a specified area. The probability construction is
\begin{equation}
p_h^{\Delta T(k+t)}
=\frac{N_h^{\Delta T(k+t)}}{\displaystyle\sum_{g=1}^{H(k+t)}N_g^{\Delta T(k+t)}}.
\tag{12}
\end{equation}
Here $N_h^{\Delta T(k+t)}$ is the total number of travel requests for origin--destination cluster pair $h$ during the specified interval, from the data used to construct the probabilities. These counts are fixed data, not the prediction horizon $N$. The source also states
\[
\sum_{h=1}^{H(k+t)}p_h^{\Delta T(k+t)}=1.
\]
Equation (15) writes these probability parameters with an additional event argument, $p_h^{\Delta T(k+t)}(k+t)$; (12) omits that argument. The argument is retained at its printed use sites. No zero-denominator convention is given for (12).

The two-step prediction-time parameter is $\tau$. Physically it represents the expected time for a predicted request to occur, but the authors emphasize its role as the best instant for inserting the future expected call to optimize the routing scheme. It is a continuous tuning parameter calibrated by sensitivity analysis, rather than a routing decision in (15). No general numerical value, bound, or equation connecting $\tau$ to $T_j^0$ is printed. The example-specific calibrated values and demand distributions are not conditions of this general formulation.

\section*{Printed objective-chain inconsistency}
For reference to the selected objective's derivation, the parent incremental-cost expression is
\begin{equation}
\operatorname*{Min}_{S(k)}J
=\sum_{t=1}^{N}\sum_{j=1}^{F}\sum_{h=1}^{H(k+t)}
p_h^{\Delta T(k+t)}
\left(\left.\bigl(C_j(k+t)-C_j(k+t-1)\bigr)
\right|_{S_j(k+t-2),h}\right),
\qquad N=2\text{ here}.
\tag{10}
\end{equation}
The original page image of (15), however, prints the following first-line expression:
\[
\sum_{t=1}^{2}\sum_{j=1}^{F}\sum_{h=1}^{H(k+t)}
p_h^{\Delta T(k+t)}(k+t)\cdot
\left.\bigl(C_j(k+t)-C_j(k+t+1)\bigr)
\right|_{S_j(k+t-2),h}.
\]
The printed $k+t+1$ contradicts both (10) and the subsequent expansion in (15). The latter expansion subtracts $C_j(k)$, cancels the first-step cost using the second-step probabilities summing to one, and ends with the selected objective at the start of this article. The conflicting first line is recorded only as a source inconsistency; it is not an additional objective or constraint.

\paragraph{Scope of the reconstruction.}
The operative formulas above are from Cort\'es, S\'aez, N\'u\~nez, and Mu\~noz-Carpintero, \emph{Hybrid Adaptive Predictive Control for a Dynamic Pickup and Delivery Problem}, Transportation Science 43(1), 27--42 (2009), especially physical PDF pages 4--8 and the insertion-feasibility discussion on pages 9--11. The supplied original pages provide (11), (12), (16), and the complete matrices in (8)--(9); these are not unavailable dependencies. Nevertheless, the source does not give a fully algebraic capacity-feasibility or event-time advancement specification for this model. No missing inequality, penalty, travel-time law, boundary convention, or transition rule has been supplied by reconstruction. No official supplement was supplied; any supplemental specification remains unverified.

\end{document}
